\documentclass[11pt,a4paper]{article}

\usepackage[T1]{fontenc}
\usepackage[utf8]{inputenc}
\usepackage{lmodern}
\usepackage{microtype}
\usepackage[margin=2.5cm]{geometry}
\usepackage{amsmath}
\usepackage{amssymb}
\usepackage{booktabs}
\usepackage{hyperref}

\hypersetup{
  colorlinks=true,
  linkcolor=blue!55!black,
  urlcolor=blue!55!black,
  citecolor=blue!55!black
}

\title{Understanding Overdispersion in the IoT Traffic Datasets\\[0.4em]
\large From the Multinomial Model to the Dirichlet--Multinomial and the LM Likelihood}
\author{}
\date{}

\begin{document}

\maketitle

\begin{abstract}
This document explains the different meanings of overdispersion that arise in
the project and in the supporting papers.  It begins with the Multinomial and
Dirichlet--Multinomial models, explains the roles of the parameters
$\boldsymbol{\alpha}$, $A$, $\psi$, and $\rho$, and then separates the
statistical question of whether the data are overdispersed from the numerical
question of how the likelihood is evaluated.  The fitted Chromecast and
Samsung Camera models are used throughout as concrete examples.  Particular
attention is paid to the criterion $\psi<0.05$ used in the ALM IoT paper,
because its direction is not the direction implied by the standard
Dirichlet--Multinomial variance formula.
\end{abstract}

\tableofcontents
\newpage

\section{Two Questions That Must Be Kept Separate}

The discussion becomes much easier once two different questions are separated.
The first is a statistical question: does the composition of the observed
traffic vary more than a fixed-probability Multinomial model predicts?  The
second is a numerical question: once a Dirichlet--Multinomial model has been
chosen, can its log-likelihood be calculated accurately and efficiently on a
computer?

The Dirichlet--Multinomial distribution addresses the statistical question.
The Yu--Shaw (YS) and Languasco--Migliardi (LM) procedures address the numerical
question.  YS and LM are not competing probability models and they do not
represent different kinds of overdispersion.  They are alternative ways of
calculating the same Dirichlet--Multinomial log-likelihood.

This distinction is important because the statistically least dispersed part
of the parameter space is also a numerically difficult part of the parameter
space.  In other words, ``close to the Multinomial model'' and ``easy to
calculate'' are not synonymous.

\section{What Is Being Modelled?}

The project divides network traffic into fixed time windows.  Each window
produces a four-category count vector

\[
\mathbf{x}_i=
\bigl(
x_{i,\mathrm{TCP}},
x_{i,\mathrm{UDP}},
x_{i,\mathrm{SSDP}},
x_{i,\mathrm{ARP}}
\bigr),
\]

whose total packet count is

\[
n_i=\sum_{k=1}^{4}x_{ik}.
\]

The Multinomial and Dirichlet--Multinomial models describe how this observed
total is divided among the four protocol categories.  They condition on
$n_i$.  They do not, by themselves, explain why one window contains 500
packets and another contains 10,000.  Their principal concern is whether those
packets are divided among TCP, UDP, SSDP, and ARP in a stable or variable way.

This point also explains why a Multinomial dispersion diagnostic is not the
same as simply calculating the variance of the total number of packets.  A
device may change its total traffic volume while retaining almost the same
protocol proportions, or it may retain a similar total while changing its
protocol mixture dramatically.  These are different behaviours.

\section{The Ordinary Multinomial Model}

The ordinary Multinomial model assumes a single probability vector

\[
\mathbf{p}=
\bigl(
p_{\mathrm{TCP}},
p_{\mathrm{UDP}},
p_{\mathrm{SSDP}},
p_{\mathrm{ARP}}
\bigr),
\qquad
\sum_k p_k=1,
\]

and uses that same vector in every window:

\[
\mathbf{X}_i\mid n_i,\mathbf{p}
\sim \operatorname{Multinomial}(n_i,\mathbf{p}).
\]

For protocol category $k$, the conditional mean and variance are

\begin{align}
\operatorname{E}[X_{ik}\mid n_i]
  &=n_i p_k,\\
\operatorname{Var}_{\mathrm{MN}}(X_{ik}\mid n_i)
  &=n_i p_k(1-p_k).
\end{align}

Suppose, for example, that UDP has probability $0.90$.  The Multinomial model
allows ordinary sampling fluctuation around 90 percent, but it assumes that
the underlying 90 percent probability remains fixed from one window to the
next.

That assumption is often too restrictive for an IoT device.  Even during
normal operation, a device may move between idle, discovery, cloud
communication, streaming, maintenance, and user-driven modes.  Each mode can
have a different protocol mixture.  If those mixtures change more than the
fixed-$\mathbf{p}$ Multinomial permits, the observations exhibit
\emph{extra-Multinomial variation}, usually called Multinomial
overdispersion.

Overdispersion is therefore always relative to a reference model.  It does not
merely mean that a raw variance looks large; it means that the observed
variation is larger than the variation predicted by the model being assessed.

\section{The Dirichlet--Multinomial Model}

The Dirichlet--Multinomial model relaxes the fixed-probability assumption.  It
allows every window to have its own protocol probability vector
$\mathbf{q}_i$.  The model can be understood as a two-stage process:

\begin{align}
\mathbf{q}_i
  &\sim \operatorname{Dirichlet}(\boldsymbol{\alpha}),\\
\mathbf{X}_i\mid n_i,\mathbf{q}_i
  &\sim \operatorname{Multinomial}(n_i,\mathbf{q}_i).
\end{align}

The contrast with the ordinary Multinomial can be summarized as

\[
\begin{array}{lll}
\text{Multinomial:}
  & \mathbf{p} & \longrightarrow \mathbf{X}_i,\\[0.4em]
\text{Dirichlet--Multinomial:}
  & \boldsymbol{\alpha} \longrightarrow \mathbf{q}_i
  & \longrightarrow \mathbf{X}_i.
\end{array}
\]

The random vector $\mathbf{q}_i$ represents the protocol mixture active in a
particular window.  One window may be more UDP-heavy, another more TCP-heavy,
and another may contain an unusual amount of ARP traffic.  The Dirichlet
distribution describes how these window-specific mixtures vary around their
long-run average.  Once the random probabilities are integrated out, the
resulting count distribution is the Dirichlet--Multinomial.

\subsection{The meaning of $\boldsymbol{\alpha}$ and $A$}

Let

\[
\boldsymbol{\alpha}=(\alpha_1,\ldots,\alpha_K),
\qquad
A=\sum_{k=1}^{K}\alpha_k.
\]

The mean probability of category $k$ is

\[
p_k=\frac{\alpha_k}{A}.
\]

The ratios among the $\alpha_k$ values therefore determine the average traffic
composition.  Their sum $A$, called the total concentration, determines how
tightly the window-specific vectors $\mathbf{q}_i$ remain around that average.
A large concentration produces similar probability vectors from window to
window.  A small concentration permits much greater changes.

This distinction can be seen by multiplying every component of
$\boldsymbol{\alpha}$ by the same positive constant.  The mean probabilities
$\alpha_k/A$ do not change, because both numerator and denominator are scaled
equally.  The concentration does change, however, so the amount of
window-to-window variability changes even though the average protocol mixture
does not.

\section{How $\psi$ and $\rho$ Measure Dispersion}

The project and the ALM papers define the reciprocal-concentration parameter

\[
\psi=\frac{1}{A}.
\]

A closely related standard Dirichlet--Multinomial parameter is

\[
\rho=\frac{1}{A+1}.
\]

Substituting $A=1/\psi$ gives a direct relationship between them:

\[
\rho=\frac{\psi}{1+\psi}.
\]

Both quantities move in the same direction.  As $\psi$ increases, $\rho$
also increases.  The importance of this fact becomes clear from the
Dirichlet--Multinomial variance:

\[
\operatorname{Var}_{\mathrm{DM}}(X_{ik}\mid n_i)
=
n_i p_k(1-p_k)
\left[1+(n_i-1)\rho\right].
\]

The first part, $n_i p_k(1-p_k)$, is exactly the Multinomial variance.  The
bracketed part is a variance-inflation factor:

\[
F_i=1+(n_i-1)\rho.
\]

The direction is consequently unambiguous.  When $\psi$ approaches zero,
$A$ approaches infinity, $\rho$ approaches zero, and $F_i$ approaches one.
The Dirichlet--Multinomial then approaches the ordinary Multinomial.  When
$\psi$ increases, the concentration decreases and the extra-Multinomial
variance increases.

\subsection{A small numerical example}

Consider a two-category example with $n=100$ and $p=0.5$.  The ordinary
Multinomial variance for the first category is

\[
100(0.5)(0.5)=25.
\]

If $\psi=0.01$, then $\rho=0.01/1.01\approx0.0099$ and

\[
F\approx1+99(0.0099)=1.98.
\]

The corresponding Dirichlet--Multinomial variance is approximately
$25\times1.98=49.5$.  If $\psi$ is increased to $0.05$, then
$\rho\approx0.0476$ and $F\approx5.71$, producing a variance of approximately
$142.9$.  The larger value of $\psi$ therefore produces the larger variance.

This example also shows why a numerically small value of $\rho$ need not have
a small practical effect.  The inflation contains the factor $n-1$.  When a
window contains thousands of packets, even modest within-window dependence
can produce substantial extra variation.

\section{The Bioinformatics Paper: A Numerical Problem Near $\psi=0$}

Yu and Shaw describe the Dirichlet--Multinomial as a model for
extra-Multinomial count variation and state that their reciprocal-concentration
parameter measures departure from the Multinomial model
\cite{yu-shaw-2014}.  Greater reciprocal concentration means greater variance;
at zero, the Dirichlet--Multinomial reduces to the Multinomial.

Their main contribution, however, is computational.  The conventional
Dirichlet--Multinomial log-likelihood contains paired differences such as

\[
\log\Gamma(A+n)-\log\Gamma(A).
\]

When $\psi$ is very small, $A=1/\psi$ is very large.  The two log-gamma terms
can then be individually enormous while their difference is comparatively
small.  In exact mathematics there is no problem, but floating-point
arithmetic stores only a finite number of significant digits.  Subtracting
two large, nearly equal approximations can discard many of those digits, a
phenomenon known as catastrophic cancellation.

Yu and Shaw introduced a series-based formulation and a mesh procedure to
evaluate the likelihood more accurately in this regime.  Their result creates
an important distinction:

\begin{quote}
Small $\psi$ means statistically close to the Multinomial boundary, but it can
also mean numerically difficult for the ordinary paired log-gamma calculation.
\end{quote}

Thus, numerical difficulty near $\psi=0$ must not be interpreted as evidence
that overdispersion becomes stronger near $\psi=0$.  These are different
ideas: one concerns the stability of a computer calculation, while the other
concerns the variance of a probability model.

\section{The Professor's Papers and the ALM Criterion}

The ALM IoT paper uses a fixed-point method to fit Dirichlet--Multinomial
parameters to aggregated IoT traffic and compares two procedures for
calculating the likelihood: Yu--Shaw and Languasco--Migliardi
\cite{alm-iot}.  Its experiments report that LM is faster and more accurate in
the tested settings.  The later TPAMI paper performs a broader comparison and
again reports advantages for LM \cite{baddar-et-al-2025}.

The crucial point for interpreting the project is that LM is a numerical
backend.  It helps the estimator evaluate the likelihood reliably and decide
whether successive parameter estimates have converged.  LM does not create
overdispersion, and using LM does not prove that a dataset is overdispersed.
In exact arithmetic, LM, YS, and a direct log-gamma calculation all target the
same Dirichlet--Multinomial likelihood.

The ALM IoT paper additionally adopts the empirical rule

\[
\psi<0.05
\]

for declaring a dataset overdispersed.  With the paper's definition
$\psi=1/A$, this inequality is equivalent to

\[
A>20
\qquad\text{and}\qquad
\rho<\frac{0.05}{1.05}\approx0.047619.
\]

It therefore selects fits with relatively large concentration and relatively
small Dirichlet--Multinomial dispersion.  This is the opposite direction from
the standard variance interpretation.  The same convention appears in the
later TPAMI paper, so it is safest to report it as an empirical classification
rule used in those experiments, not as a general theorem or formal
hypothesis test.  In particular, the number $0.05$ is not a $p$-value or a
statistical significance level in this context.

One possible source of the confusing language is that the papers focus on the
numerically delicate low-$\psi$ regime.  Low $\psi$ can still produce
substantial variance inflation when each observation has a very large total
count.  Neither point changes the monotonic relationship in the variance
formula: for the same $n$ and $\mathbf{p}$, increasing $\psi$ increases the
Dirichlet--Multinomial variance.

\section{The Chromecast and Samsung Training Models}

The fitted parameters stored by the project are shown in
Table~\ref{tab:training-models}.  The final column asks a simple hypothetical
question: if both models were applied to a window with $n=100$, by what factor
would the Dirichlet--Multinomial variance exceed the corresponding
Multinomial variance?

\begin{table}[htbp]
\centering
\caption{Fitted training dispersion parameters.}
\label{tab:training-models}
\begin{tabular}{lrrrr}
\toprule
Dataset & $A$ & $\psi=1/A$ & $\rho=\psi/(1+\psi)$ & $F$ at $n=100$ \\
\midrule
Chromecast & 28.0985 & 0.035589 & 0.034366 & 4.40 \\
Samsung Camera & 17.3941 & 0.057491 & 0.054365 & 6.38 \\
\bottomrule
\end{tabular}
\end{table}

At the same window total, the fitted Samsung model is the more dispersed of
the two.  Its concentration is lower, its $\psi$ is higher, its $\rho$ is
higher, and consequently its variance-inflation factor is higher.

The fitted mean protocol probabilities provide different information.  For
Chromecast, the fitted normal mixture is approximately 92.56 percent UDP,
5.42 percent TCP, 2.02 percent ARP, and 0.004 percent SSDP.  For Samsung, it is
approximately 51.00 percent SSDP, 36.57 percent TCP, 12.26 percent ARP, and
0.17 percent UDP.  These percentages describe the centres of the two fitted
distributions.  They should not be confused with concentration or
overdispersion, which describe how much individual windows move around those
centres.

Under the ALM paper's empirical rule, Chromecast ``passes'' because
$0.035589<0.05$, while Samsung ``fails'' because $0.057491>0.05$.  This wording
does not mean that Samsung lacks overdispersion.  It means only that Samsung
does not satisfy the particular inequality adopted by that paper.  Indeed,
under the standard Dirichlet--Multinomial interpretation, Samsung has the
larger fitted dispersion parameter.

The model values are stored in
\path{artifacts/UNSW-Chromecast/model.json} and
\path{artifacts/UNSW-Samsung-Camera/model.json}.

\section{Why the Pearson Result Has a Different Ordering}

The direct Pearson Multinomial diagnostic produced the training values

\[
\phi_{\mathrm{Chromecast}}=218.96,
\qquad
\phi_{\mathrm{Samsung}}=86.84.
\]

A conventional Pearson dispersion ratio has the form

\[
\phi=\frac{X^2}{\text{residual degrees of freedom}},
\qquad
X^2=
\sum_{i,k}
\frac{(x_{ik}-n_i\widehat{p}_k)^2}
     {n_i\widehat{p}_k}.
\]

A value near one is expected when the fixed-probability Multinomial adequately
describes the residual variation.  Values much greater than one indicate
extra-Multinomial variation.  Both training results are therefore strong
evidence that a fixed-$\mathbf{p}$ Multinomial is too narrow for these data.

It may initially seem contradictory that Samsung has the larger fitted
$\psi$, while Chromecast has the larger Pearson $\phi$.  There is no
contradiction because the two quantities are not the same statistic.  In
particular, the Dirichlet--Multinomial inflation factor depends on both
$\rho$ and the total number of packets in a window:

\[
F_i=1+(n_i-1)\rho.
\]

The mean training-window totals are approximately 11,481 packets for
Chromecast and 615 packets for Samsung.  Substituting these mean totals merely
as an illustration gives

\begin{align}
F_{\mathrm{Chromecast}}
  &\approx1+(11{,}481-1)(0.034366)
    \approx395.5,\\
F_{\mathrm{Samsung}}
  &\approx1+(615-1)(0.054365)
    \approx34.4.
\end{align}

Samsung has greater fitted probability heterogeneity when $n$ is held fixed,
but Chromecast windows contain far more packets.  Even a smaller
within-window correlation can generate very large extra variation when it is
shared by thousands of packet outcomes.

These illustrative inflation factors are not expected to equal the observed
Pearson ratios.  The real window totals vary, the probabilities are estimated,
some categories are sparse, and one stationary Dirichlet--Multinomial is only
an approximation to traffic that changes over time.  The calculation is
useful for explaining the direction and scale, not for predicting the Pearson
statistic exactly.

\section{What the Individual Days Reveal}

The daily fits make the direction of $\psi$ particularly easy to see.  For
Chromecast, October 10 has $\psi=0.016012$, whereas October 14 has
$\psi=0.072661$.  Under the standard model interpretation, October 14 is the
more dispersed day.  Nevertheless, the ALM inequality labels October 10 a
pass and October 14 a fail.  October 23, with $\psi=0.196667$, has the greatest
fitted dispersion among the listed Chromecast days, not the smallest.

The Samsung data show an even sharper contrast.  May 29 has
$\psi=0.000942$ and lies very close to the Multinomial boundary.  May 28 has
$\psi=0.147219$ and allows much greater compositional variation.  The ALM rule
again calls the first a pass and the second a fail because its inequality runs
in the direction of smaller $\psi$.

The pooled Samsung training value, $\psi=0.057491$, is not the arithmetic
average of its daily values.  A pooled fit estimates one model from all
windows simultaneously.  It is nonlinear, gives different effective weight
to different observations, and also sees changes between daily traffic
regimes.  In particular, the highly variable May 28 traffic and the shift from
May 28 to the much more stable May 29 and May 30 traffic contribute to the
heterogeneity of the pooled fit.

The two completely silent Chromecast days have no daily Dirichlet--Multinomial
fit.  When every window has total count zero, there is no observed protocol
composition from which to estimate the category probabilities or their
dispersion.

Daily results are recorded in
\path{reports/UNSW-Chromecast/capture_statistics.json} and
\path{reports/UNSW-Samsung-Camera/capture_statistics.json}.

\section{Why Pooling Partitions Changes the Result}

Table~\ref{tab:aggregate} shows that the fitted value of $\psi$ changes
substantially across data partitions.  The labels in the final column apply
only the ALM paper's empirical inequality and should not be read as a standard
test of overdispersion.

\begin{table}[htbp]
\centering
\caption{Aggregate fitted values and the ALM paper's empirical classification.}
\label{tab:aggregate}
\begin{tabular}{llrc}
\toprule
Dataset & Scope & $\psi$ & ALM rule $\psi<0.05$ \\
\midrule
Chromecast & Fit/training & 0.035589 & Pass \\
Chromecast & Calibration & 0.058771 & Fail \\
Chromecast & Development & 0.119722 & Fail \\
Chromecast & Final & 0.032168 & Pass \\
Chromecast & Entire dataset & 0.076318 & Fail \\
\addlinespace
Samsung Camera & Fit/training & 0.057491 & Fail \\
Samsung Camera & Calibration & 0.008272 & Pass \\
Samsung Camera & Development & 0.127651 & Fail \\
Samsung Camera & Final & 0.038767 & Pass \\
Samsung Camera & Entire dataset & 0.073784 & Fail \\
\bottomrule
\end{tabular}
\end{table}

Pooling different days asks one stationary model to explain all their windows.
If normal behaviour changes with time, or if normal and attack periods are
combined, the pooled model interprets some of those regime changes as
additional compositional heterogeneity.  A pooled $\psi$ can therefore be
larger than a training $\psi$ even when several individual days have small
values.

This is also why a large whole-dataset value does not prove that one stationary
Dirichlet--Multinomial fits every day well.  It may partly reflect genuine
within-regime variability, but it can also reflect unmodelled changes between
regimes.  In an anomaly-detection setting, attacks and temporal changes should
not automatically be absorbed into the definition of normal dispersion.

\section{Three Diagnostics That Should Not Be Confused}

The project reports several quantities that all contain the word
``dispersion'', but they answer different questions.

First, a per-category variance-to-mean index,

\[
\frac{\operatorname{Var}(X_k)}{\operatorname{E}[X_k]},
\]

is a Poisson-style diagnostic for one category.  It is affected by changes in
total traffic volume and is not, by itself, a direct measure of Multinomial
compositional overdispersion.

Second, the Pearson ratio $\phi$ compares complete count vectors with the
variation expected from a fixed-probability Multinomial, while conditioning on
each window's total.  It is a direct empirical diagnostic of
extra-Multinomial variation.  Values far above one motivate a richer model.

Third, $\psi=1/A$ is a fitted Dirichlet--Multinomial parameter.  It describes
the dispersion inside that fitted model: larger $\psi$ means smaller
concentration and greater variance for fixed $n$ and $\mathbf{p}$.  It is not
itself a goodness-of-fit test.  Since the fitting procedure constrains the
Dirichlet parameters to be positive, obtaining a positive $\psi$ alone is not
proof that the Dirichlet--Multinomial is necessary.  That case is instead
supported by diagnostics such as Pearson $\phi$, goodness-of-fit comparisons,
or explicit comparison with the simpler Multinomial model.

Finally, the statement ``LM was used'' says only how the likelihood was
calculated.  It is neither a dispersion statistic nor a test result.

\section{A Careful Interpretation for the Thesis}

The statistical evidence supports the following interpretation:

\begin{quote}
The normal training partitions exhibit strong extra-Multinomial variation, as
shown by Pearson dispersion ratios substantially greater than one.  This
supports using a Dirichlet--Multinomial rather than a fixed-probability
Multinomial.  The fitted reciprocal-concentration parameter is
$\psi=0.035589$ for Chromecast and $\psi=0.057491$ for Samsung Camera.  Under
the standard Dirichlet--Multinomial variance interpretation, the larger
Samsung value represents greater probability heterogeneity when the window
total is held fixed.  Chromecast satisfies the empirical $\psi<0.05$
convention adopted by the ALM paper, whereas Samsung does not.  This
convention is reported for comparison only, because its inequality runs in the
opposite direction from the usual interpretation of $\psi=1/A$.  The LM
procedure is used as an accurate numerical implementation of the
Dirichlet--Multinomial log-likelihood, not as a test for overdispersion.
\end{quote}

The shortest conceptual summary is therefore:

\[
\boxed{
\begin{gathered}
\text{larger }\psi
\Longleftrightarrow
\text{smaller }A
\Longleftrightarrow
\text{larger }\rho
\Longleftrightarrow
\text{greater DM variance},\\[0.4em]
\psi\longrightarrow0
\Longleftrightarrow
\text{the DM approaches the Multinomial, while ordinary numerical evaluation
becomes delicate.}
\end{gathered}}
\]

\begin{thebibliography}{9}

\bibitem{yu-shaw-2014}
P. Yu and C. A. Shaw,
``An efficient algorithm for accurate computation of the
Dirichlet-multinomial log-likelihood function,''
\emph{Bioinformatics}, vol. 30, no. 11, pp. 1547--1554, 2014.
\url{https://doi.org/10.1093/bioinformatics/btu079}.

\bibitem{alm-iot}
S. Al-Haj Baddar, A. Languasco, and M. Migliardi,
``Modeling and forecasting overdispersed IoT data using an efficient and
accurate computation of the Dirichlet Multinomial distribution.''
Local copy: \path{resouces/ALM_IoT-CR.pdf}.

\bibitem{baddar-et-al-2025}
S. Al-Haj Baddar, A. Languasco, and M. Migliardi,
``Efficient analysis of overdispersed data using an accurate computation of
the Dirichlet Multinomial distribution,''
\emph{IEEE Transactions on Pattern Analysis and Machine Intelligence},
vol. 47, no. 2, pp. 1181--1189, 2025.
\url{https://doi.org/10.1109/TPAMI.2024.3489645}.

\end{thebibliography}

\end{document}
